%% Generated by Sphinx.
\def\sphinxdocclass{report}
\documentclass[letterpaper,10pt,english]{sphinxmanual}
\ifdefined\pdfpxdimen
   \let\sphinxpxdimen\pdfpxdimen\else\newdimen\sphinxpxdimen
\fi \sphinxpxdimen=.75bp\relax
\ifdefined\pdfimageresolution
    \pdfimageresolution= \numexpr \dimexpr1in\relax/\sphinxpxdimen\relax
\fi
%% let collapsible pdf bookmarks panel have high depth per default
\PassOptionsToPackage{bookmarksdepth=5}{hyperref}

\PassOptionsToPackage{booktabs}{sphinx}
\PassOptionsToPackage{colorrows}{sphinx}

\PassOptionsToPackage{warn}{textcomp}
\usepackage[utf8]{inputenc}
\ifdefined\DeclareUnicodeCharacter
% support both utf8 and utf8x syntaxes
  \ifdefined\DeclareUnicodeCharacterAsOptional
    \def\sphinxDUC#1{\DeclareUnicodeCharacter{"#1}}
  \else
    \let\sphinxDUC\DeclareUnicodeCharacter
  \fi
  \sphinxDUC{00A0}{\nobreakspace}
  \sphinxDUC{2500}{\sphinxunichar{2500}}
  \sphinxDUC{2502}{\sphinxunichar{2502}}
  \sphinxDUC{2514}{\sphinxunichar{2514}}
  \sphinxDUC{251C}{\sphinxunichar{251C}}
  \sphinxDUC{2572}{\textbackslash}
\fi
\usepackage{cmap}
\usepackage[T1]{fontenc}
\usepackage{amsmath,amssymb,amstext}
\usepackage{babel}



\usepackage{tgtermes}
\usepackage{tgheros}
\renewcommand{\ttdefault}{txtt}



\usepackage[Bjarne]{fncychap}
\usepackage{sphinx}

\fvset{fontsize=auto}
\usepackage{geometry}


% Include hyperref last.
\usepackage{hyperref}
% Fix anchor placement for figures with captions.
\usepackage{hypcap}% it must be loaded after hyperref.
% Set up styles of URL: it should be placed after hyperref.
\urlstyle{same}

\addto\captionsenglish{\renewcommand{\contentsname}{Spis Treści:}}

\usepackage{sphinxmessages}
\setcounter{tocdepth}{1}



\title{Sprawozdanie\sphinxhyphen{}z\sphinxhyphen{}laboratoriów}
\date{May 19, 2026}
\release{}
\author{Paweł Łoćwin}
\newcommand{\sphinxlogo}{\vbox{}}
\renewcommand{\releasename}{}
\makeindex
\begin{document}

\ifdefined\shorthandoff
  \ifnum\catcode`\=\string=\active\shorthandoff{=}\fi
  \ifnum\catcode`\"=\active\shorthandoff{"}\fi
\fi

\pagestyle{empty}
\sphinxmaketitle
\pagestyle{plain}
\sphinxtableofcontents
\pagestyle{normal}
\phantomsection\label{\detokenize{index::doc}}


\sphinxAtStartPar
\sphinxstylestrong{Autorzy projektu:} Paweł Łoćwin Paweł Łosowski

\sphinxAtStartPar
Niniejsza dokumentacja stanowi kompletne sprawozdanie z laboratorium przedmiotu Bazy Danych. Projekt kompleksowo obejmuje aspekty związane z zarządzaniem środowiskiem bazodanowym, projektowaniem struktury relacyjnej oraz implementacją fizyczną.

\sphinxAtStartPar
Dokumentacja została podzielona na cztery główne sekcje tematyczne, począwszy od teoretycznych aspektów konserwacji systemów, poprzez badania literaturowe, aż po praktyczny projekt i wdrożenie bazy danych dla systemu zarządzania biblioteką.

\sphinxstepscope


\chapter{Rozdział 1: Kontrola i konserwacja}
\label{\detokenize{rozdzial_1/index:rozdzial-1-kontrola-i-konserwacja}}\label{\detokenize{rozdzial_1/index::doc}}\begin{quote}\begin{description}
\sphinxlineitem{Autorzy}\begin{enumerate}
\sphinxsetlistlabels{\arabic}{enumi}{enumii}{}{.}%
\item {} 
\sphinxAtStartPar
Paweł Łoćwin

\item {} 
\sphinxAtStartPar
Paweł Łosowski

\end{enumerate}

\end{description}\end{quote}


\section{Wprowadzenie}
\label{\detokenize{rozdzial_1/index:wprowadzenie}}
\sphinxAtStartPar
Współczesne systemy relacyjnych baz danych to wysoce skomplikowane środowiska, które wymagają ciągłego i proaktywnego nadzoru, aby zagwarantować wysoką dostępność, niezawodność oraz optymalny czas odpowiedzi na zapytania klientów. Wraz z rosnącym wolumenem przetwarzanych informacji, zarządzanie cyklem życia danych staje się wyzwaniem równie ważnym, co sam projekt struktury relacyjnej. Rozdział ten poświęcony jest nowoczesnym mechanizmom kontroli i konserwacji, ze szczególnym uwzględnieniem architektury silnika PostgreSQL. Procesy te obejmują zaawansowane zarządzanie przestrzenią dyskową, analizę i optymalizację planów zapytań, rygorystyczną kontrolę dostępu oraz wielowarstwowe strategie zabezpieczania danych przed ich bezpowrotną utratą.


\section{Zarządzanie przestrzenią i mechanizm współbieżności}
\label{\detokenize{rozdzial_1/index:zarzadzanie-przestrzenia-i-mechanizm-wspolbieznosci}}
\sphinxAtStartPar
PostgreSQL opiera swoje działanie na zaawansowanej architekturze określanej jako Multi\sphinxhyphen{}Version Concurrency Control. Technologia ta pozwala na równoległy odczyt i zapis danych bez wzajemnego blokowania się transakcji, co drastycznie zwiększa przepustowość systemu przy dużej liczbie jednoczesnych połączeń. Skutkiem ubocznym tego bezkolizyjnego mechanizmu jest jednak powstawanie martwych krotek, czyli przestarzałych wersji wierszy. Wiersze te, choć zostały już zaktualizowane lub usunięte przez aplikację, wciąż fizycznie zajmują przestrzeń na dysku ze względu na to, że inne, wciąż otwarte transakcje, mogą potrzebować dostępu do ich historycznego stanu.

\sphinxAtStartPar
Aby zapobiec zjawisku drastycznego puchnięcia tabel oraz degradacji wydajności operacji wejścia i wyjścia, konieczna jest regularna konserwacja na poziomie nośnika danych:
\begin{itemize}
\item {} 
\sphinxAtStartPar
\sphinxstylestrong{VACUUM:} Podstawowy proces konserwacyjny odzyskujący miejsce po martwych krotkach. Oznacza on wolne obszary jako dostępne do ponownego zapisu przez nowe instrukcje, zapobiegając niekontrolowanemu rozrostowi plików bazy.

\item {} 
\sphinxAtStartPar
\sphinxstylestrong{VACUUM FULL:} Znacznie bardziej inwazyjna wersja operacji, która fizycznie przebudowuje strukturę całej tabeli i trwale zwraca wolne miejsce do systemu operacyjnego. Operacja ta wymaga jednak nałożenia wyłącznej blokady na modyfikowany obiekt, co uniemożliwia jego odczyt i zapis na czas trwania przebudowy.

\item {} 
\sphinxAtStartPar
\sphinxstylestrong{AUTOVACUUM:} W nowoczesnych środowiskach produkcyjnych utrzymanie czystości dyskowej powierza się wbudowanemu procesowi pracującemu w tle. Monitoruje on na bieżąco statystyki zmian i automatycznie inicjuje procesy czyszczenia po przekroczeniu zdefiniowanych progów aktywności, co minimalizuje konieczność ręcznej ingerencji administratora.

\end{itemize}

\sphinxAtStartPar
Zaniedbanie procesu czyszczenia w systemach o wysokiej rotacji danych może doprowadzić nie tylko do spadku wydajności, ale w skrajnych przypadkach do wyczerpania identyfikatorów transakcji, co skutkuje całkowitym wstrzymaniem działania bazy.

\begin{sphinxVerbatim}[commandchars=\\\{\}]
\PYG{c+c1}{\PYGZhy{}\PYGZhy{} Przykładowe wymuszenie manualnej konserwacji z analizą statystyk dla tabeli}
\PYG{k}{VACUUM}\PYG{+w}{ }\PYG{p}{(}\PYG{k}{VERBOSE}\PYG{p}{,}\PYG{+w}{ }\PYG{k}{ANALYZE}\PYG{p}{)}\PYG{+w}{ }\PYG{n}{Wypozyczenia}\PYG{p}{;}
\end{sphinxVerbatim}


\section{Optymalizacja i statystyki planisty zapytań}
\label{\detokenize{rozdzial_1/index:optymalizacja-i-statystyki-planisty-zapytan}}
\sphinxAtStartPar
Kolejnym filarem utrzymania sprawności systemu jest kontrola nad optymalizatorem, zwanym planistą zapytań. Silnik bazy danych przed wykonaniem jakiejkolwiek skomplikowanej kwerendy analizuje dziesiątki możliwych ścieżek jej realizacji. Decyzje takie jak wybór między skanowaniem sekwencyjnym całego zbioru a użyciem konkretnego indeksu podejmowane są na podstawie modelu matematycznego opartego o szacowany koszt operacji.

\sphinxAtStartPar
Aby planista mógł podejmować optymalne decyzje, musi opierać się na wysoce dokładnych i aktualnych statystykach dotyczących rozkładu danych. Obejmują one między innymi histogramy wartości, listy najczęściej występujących elementów oraz współczynniki korelacji między kolumnami.
\begin{itemize}
\item {} 
\sphinxAtStartPar
Instrukcja \sphinxstylestrong{ANALYZE} służy do natychmiastowego odświeżenia tych metadanych. Baza pobiera losową próbkę wierszy i na jej podstawie aktualizuje wewnętrzne tabele systemowe, co pozwala uniknąć katastrofalnych w skutkach błędów estymacji optymalizatora.

\item {} 
\sphinxAtStartPar
Bieżące profilowanie wydajności realizowane jest natomiast przy pomocy instrukcji \sphinxstylestrong{EXPLAIN ANALYZE}. Narzędzie to uruchamia zapytanie, po czym zwraca nie tylko wynik, ale również szczegółowy raport z rzeczywistego przebiegu wykonania. Pozwala to programistom i administratorom precyzyjnie namierzyć wąskie gardła w skomplikowanych kwerendach, które łączą wiele tabel i filtrują ogromne zbiory danych.

\end{itemize}

\begin{sphinxVerbatim}[commandchars=\\\{\}]
\PYG{c+c1}{\PYGZhy{}\PYGZhy{} Kontrola planu wykonania zapytania z rzeczywistym pomiarem czasu}
\PYG{k}{EXPLAIN}\PYG{+w}{ }\PYG{k}{ANALYZE}
\PYG{k}{SELECT}\PYG{+w}{ }\PYG{k}{c}\PYG{p}{.}\PYG{n}{Imie}\PYG{p}{,}\PYG{+w}{ }\PYG{k}{c}\PYG{p}{.}\PYG{n}{Nazwisko}\PYG{p}{,}\PYG{+w}{ }\PYG{n}{k}\PYG{p}{.}\PYG{n}{Tytul}
\PYG{k}{FROM}\PYG{+w}{ }\PYG{n}{Czytelnicy}\PYG{+w}{ }\PYG{k}{c}
\PYG{k}{JOIN}\PYG{+w}{ }\PYG{n}{Wypozyczenia}\PYG{+w}{ }\PYG{n}{w}\PYG{+w}{ }\PYG{k}{ON}\PYG{+w}{ }\PYG{k}{c}\PYG{p}{.}\PYG{n}{ID\PYGZus{}Czytelnika}\PYG{+w}{ }\PYG{o}{=}\PYG{+w}{ }\PYG{n}{w}\PYG{p}{.}\PYG{n}{ID\PYGZus{}Czytelnika}
\PYG{k}{JOIN}\PYG{+w}{ }\PYG{n}{Ksiazki}\PYG{+w}{ }\PYG{n}{k}\PYG{+w}{ }\PYG{k}{ON}\PYG{+w}{ }\PYG{n}{w}\PYG{p}{.}\PYG{n}{ID\PYGZus{}Ksiazki}\PYG{+w}{ }\PYG{o}{=}\PYG{+w}{ }\PYG{n}{k}\PYG{p}{.}\PYG{n}{ID\PYGZus{}Ksiazki}
\PYG{k}{WHERE}\PYG{+w}{ }\PYG{n}{w}\PYG{p}{.}\PYG{n}{Data\PYGZus{}Zwrotu}\PYG{+w}{ }\PYG{k}{IS}\PYG{+w}{ }\PYG{k}{NULL}\PYG{p}{;}
\end{sphinxVerbatim}


\section{Bezpieczeństwo i rygorystyczna kontrola dostępu}
\label{\detokenize{rozdzial_1/index:bezpieczenstwo-i-rygorystyczna-kontrola-dostepu}}
\sphinxAtStartPar
Bezpieczeństwo danych stanowi absolutny priorytet każdej administracji systemami informatycznymi. Nowoczesne podejście do ochrony baz danych całkowicie odchodzi od wykorzystywania globalnych kont administracyjnych, takich jak domyślny użytkownik instalacyjny, do codziennych operacji na poziomie aplikacji. Zamiast tego wdraża się model kontroli dostępu opartej na rolach, który znacząco ogranicza ryzyko kompromitacji całego systemu w przypadku wycieku pojedynczego poświadczenia.

\sphinxAtStartPar
System zabezpieczeń środowiska PostgreSQL pozwala na niezwykle elastyczną gradację uprawnień. Definiuje się dedykowane role systemowe przypisane do konkretnych mikrousług lub modułów aplikacyjnych. Poświadczenia dostępowe szyfrowane są przy wykorzystaniu najnowocześniejszych standardów kryptograficznych, takich jak algorytm SCRAM\sphinxhyphen{}SHA\sphinxhyphen{}256, odporny na ataki słownikowe i nasłuchiwanie ruchu sieciowego.

\sphinxAtStartPar
Kluczowym konceptem jest wdrożenie zasady najmniejszych uprawnień. Rola wykorzystywana przez aplikację powinna dysponować prawami pozwalającymi wyłącznie na modyfikację danych niezbędnych do jej funkcjonowania. W praktyce oznacza to celowe przyznawanie praw do wykonywania selekcji czy dodawania rekordów, przy jednoczesnym blokowaniu możliwości ich usuwania lub modyfikacji struktury tabel.

\begin{sphinxVerbatim}[commandchars=\\\{\}]
\PYG{c+c1}{\PYGZhy{}\PYGZhy{} Przykład implementacji bezpiecznej polityki kontroli dostępu}
\PYG{k}{CREATE}\PYG{+w}{ }\PYG{k}{ROLE}\PYG{+w}{ }\PYG{n}{app\PYGZus{}user}\PYG{+w}{ }\PYG{k}{WITH}\PYG{+w}{ }\PYG{n}{LOGIN}\PYG{+w}{ }\PYG{n}{PASSWORD}\PYG{+w}{ }\PYG{l+s+s1}{\PYGZsq{}TrudneHaslo123\PYGZsq{}}\PYG{p}{;}
\PYG{k}{GRANT}\PYG{+w}{ }\PYG{k}{CONNECT}\PYG{+w}{ }\PYG{k}{ON}\PYG{+w}{ }\PYG{k}{DATABASE}\PYG{+w}{ }\PYG{n}{biblioteka\PYGZus{}db}\PYG{+w}{ }\PYG{k}{TO}\PYG{+w}{ }\PYG{n}{app\PYGZus{}user}\PYG{p}{;}
\PYG{k}{GRANT}\PYG{+w}{ }\PYG{k}{USAGE}\PYG{+w}{ }\PYG{k}{ON}\PYG{+w}{ }\PYG{k}{SCHEMA}\PYG{+w}{ }\PYG{k}{public}\PYG{+w}{ }\PYG{k}{TO}\PYG{+w}{ }\PYG{n}{app\PYGZus{}user}\PYG{p}{;}
\PYG{k}{GRANT}\PYG{+w}{ }\PYG{k}{SELECT}\PYG{p}{,}\PYG{+w}{ }\PYG{k}{INSERT}\PYG{+w}{ }\PYG{k}{ON}\PYG{+w}{ }\PYG{n}{Wypozyczenia}\PYG{+w}{ }\PYG{k}{TO}\PYG{+w}{ }\PYG{n}{app\PYGZus{}user}\PYG{p}{;}
\end{sphinxVerbatim}


\section{Strategie zabezpieczania przed utratą danych}
\label{\detokenize{rozdzial_1/index:strategie-zabezpieczania-przed-utrata-danych}}
\sphinxAtStartPar
Nawet najlepiej zoptymalizowany i zabezpieczony system informatyczny jest narażony na błędy ludzkie lub awarie infrastruktury sprzętowej. Dlatego odpowiednia polityka wykonywania kopii zapasowych oraz odtwarzania środowiska po awarii jest krytycznym elementem konserwacji relacyjnych baz danych. Zarządzanie tym procesem opiera się na dwóch uzupełniających się podejściach.

\sphinxAtStartPar
Kopie logiczne polegają na ekstrakcji danych z bazy do postaci czystego skryptu języka SQL lub dedykowanego formatu archiwalnego. Narzędzia realizujące to zadanie gwarantują przenośność danych pomiędzy różnymi wersjami silnika bazy czy architekturami sprzętowymi. Proces tworzenia kopii logicznej potrafi być jednak bardzo obciążający dla procesora, a czas odtwarzania z takiego pliku w przypadku wieloterabajtowych struktur jest często nieakceptowalnie długi z biznesowego punktu widzenia.

\sphinxAtStartPar
Dlatego w krytycznych środowiskach produkcyjnych stosuje się kopie fizyczne połączone z archiwizacją dzienników wyprzedzających. Dzienniki te rejestrują absolutnie każdą zmianę bajtu na dysku jeszcze przed jej ostatecznym zatwierdzeniem. Odpowiednio skonfigurowany mechanizm ciągłej archiwizacji tych strumieni pozwala na odtworzenie całej bazy danych z chirurgiczną precyzją, do konkretnej sekundy w przeszłości. Taka funkcjonalność pozwala cofnąć skutki przypadkowego usunięcia danych bez utraty tysięcy transakcji, które miały miejsce od czasu wykonania ostatniej pełnej kopii zapasowej.


\section{Podsumowanie}
\label{\detokenize{rozdzial_1/index:podsumowanie}}
\sphinxAtStartPar
Prawidłowa kontrola i konserwacja środowiska bazodanowego to wielowymiarowy proces, który nie ma punktu końcowego, lecz trwa przez cały cykl życia oprogramowania. Złożoność nowoczesnych systemów zarządzania bazami danych wymaga od administratorów i projektantów proaktywnego, a nie jedynie reaktywnego podejścia do pojawiających się problemów.

\sphinxAtStartPar
Samo zaprojektowanie zgodnego ze sztuką schematu relacyjnego jest zaledwie początkiem pracy. Stabilność aplikacji zależy w równej mierze od rygorystycznego zarządzania fizyczną strukturą nośników danych, dogłębnego zrozumienia statystycznych mechanizmów planisty zapytań, a także bezkompromisowego wdrażania zaawansowanych polityk bezpieczeństwa opartego na podziale ról.

\sphinxAtStartPar
Holistyczne spojrzenie na administrację PostgreSQL, obejmujące zarówno automatyzację procesów porządkujących pamięć, jak i ciągłe monitorowanie planów wykonania zapytań, pozwala na budowanie rozwiązań prawdziwie skalowalnych. Dzięki zaimplementowaniu odpowiednich strategii prewencyjnych i solidnych mechanizmów odtwarzania po awarii, organizacja zyskuje pewność, że jej kluczowy zasób cyfrowy pozostanie bezpieczny, spójny i gotowy na rosnące obciążenia bez względu na niespodziewane incydenty technologiczne.

\sphinxstepscope


\chapter{Badania literaturowe}
\label{\detokenize{rozdzial_2/index:badania-literaturowe}}\label{\detokenize{rozdzial_2/index::doc}}
\sphinxstepscope


\section{Sprzęt dla bazy danych}
\label{\detokenize{rozdzial_2/rozdzial_21/index:sprzet-dla-bazy-danych}}\label{\detokenize{rozdzial_2/rozdzial_21/index::doc}}
\sphinxstepscope


\section{Konfiguracja bazy danych PostgreSQL}
\label{\detokenize{rozdzial_2/rozdzial_22/index:konfiguracja-bazy-danych-postgresql}}\label{\detokenize{rozdzial_2/rozdzial_22/index::doc}}

\subsection{Wstęp}
\label{\detokenize{rozdzial_2/rozdzial_22/index:wstep}}
\sphinxAtStartPar
Rozdział opisuje konfigurację połączenia z bazą danych \sphinxstylestrong{PostgreSQL} \textendash{} dojrzałym, open\sphinxhyphen{}source’owym systemem zarządzania relacyjnymi bazami danych (RDBMS), który wyróżnia się pełną zgodnością ze standardem SQL oraz silnym naciskiem na rozszerzalność i niezawodność.

\sphinxAtStartPar
\sphinxstylestrong{Wymagania:}
\begin{itemize}
\item {} 
\sphinxAtStartPar
dostęp do instancji PostgreSQL,

\item {} 
\sphinxAtStartPar
zmienne środowiskowe,

\item {} 
\sphinxAtStartPar
narzędzia projektowe (np. \sphinxcode{\sphinxupquote{psql}}, \sphinxcode{\sphinxupquote{pg\_isready}}).

\end{itemize}


\subsection{Podstawowe parametry}
\label{\detokenize{rozdzial_2/rozdzial_22/index:podstawowe-parametry}}
\sphinxAtStartPar
Każde połączenie z PostgreSQL wymaga:
\begin{itemize}
\item {} 
\sphinxAtStartPar
hosta i portu (domyślnie \sphinxcode{\sphinxupquote{localhost}} i \sphinxcode{\sphinxupquote{5432}}),

\item {} 
\sphinxAtStartPar
nazwy bazy danych,

\item {} 
\sphinxAtStartPar
użytkownika i hasła.

\end{itemize}

\sphinxAtStartPar
Dane wrażliwe (np. hasła) przechowuj w \sphinxstylestrong{zmiennych środowiskowych} \textendash{} to oddziela konfigurację od kodu.
Pamiętaj: zmienne środowiskowe nie są mechanizmem bezpieczeństwa; w produkcji uzupełnij je o dedykowane
zarządzanie sekretami (np. HashiCorp Vault, AWS Secrets Manager).


\subsection{Lokalizacja i struktura katalogów PostgreSQL}
\label{\detokenize{rozdzial_2/rozdzial_22/index:lokalizacja-i-struktura-katalogow-postgresql}}
\sphinxAtStartPar
Pliki konfiguracyjne PostgreSQL znajdują się w katalogu danych (\sphinxcode{\sphinxupquote{PGDATA}}).

\sphinxAtStartPar
\sphinxstylestrong{Główne pliki konfiguracyjne:}
\begin{itemize}
\item {} 
\sphinxAtStartPar
\sphinxcode{\sphinxupquote{postgresql.conf}} \textendash{} podstawowa konfiguracja serwera (port, pamięć, logi)

\item {} 
\sphinxAtStartPar
\sphinxcode{\sphinxupquote{pg\_hba.conf}} \textendash{} kontrola dostępu (kto i skąd może się łączyć)

\item {} 
\sphinxAtStartPar
\sphinxcode{\sphinxupquote{pg\_ident.conf}} \textendash{} mapowanie użytkowników systemowych na role PostgreSQL

\end{itemize}

\sphinxAtStartPar
\sphinxstylestrong{Domyślne lokalizacje:}

\begin{DUlineblock}{0em}
\item[] System | Katalog danych |
\end{DUlineblock}

\sphinxAtStartPar
{\color{red}\bfseries{}|\sphinxhyphen{}\sphinxhyphen{}\sphinxhyphen{}\sphinxhyphen{}\sphinxhyphen{}\sphinxhyphen{}\sphinxhyphen{}\sphinxhyphen{}|}—————\sphinxhyphen{}|
| Linux (RPM/DEB) | \sphinxcode{\sphinxupquote{/var/lib/pgsql/data/}} lub \sphinxcode{\sphinxupquote{/var/lib/postgresql/*/main/}} |
| Linux (kompilacja źródłowa) | \sphinxcode{\sphinxupquote{/usr/local/pgsql/data/}} |
| Windows | \sphinxcode{\sphinxupquote{C:\textbackslash{}Program Files\textbackslash{}PostgreSQL\textbackslash{}\textless{}version\textgreater{}\textbackslash{}data\textbackslash{}}} |
| Docker | \sphinxcode{\sphinxupquote{/var/lib/postgresql/data/}} (w kontenerze) |

\sphinxAtStartPar
\sphinxstylestrong{Zmiana katalogu danych:}

\begin{sphinxVerbatim}[commandchars=\\\{\}]
\PYG{c+c1}{\PYGZsh{} Inicjalizacja nowego katalogu}
initdb\PYG{+w}{ }\PYGZhy{}D\PYG{+w}{ }/ścieżka/do/katalogu

\PYG{c+c1}{\PYGZsh{} Uruchomienie z niestandardowym PGDATA}
pg\PYGZus{}ctl\PYG{+w}{ }\PYGZhy{}D\PYG{+w}{ }/ścieżka/do/katalogu\PYG{+w}{ }start
\end{sphinxVerbatim}

\sphinxAtStartPar
\sphinxstylestrong{Struktura katalogu danych:}

\begin{sphinxVerbatim}[commandchars=\\\{\}]
\PYGZdl{}PGDATA/
├── postgresql.conf   \PYGZsh{} główny plik konfiguracyjny
├── pg\PYGZus{}hba.conf       \PYGZsh{} polityka dostępu
├── pg\PYGZus{}ident.conf     \PYGZsh{} mapowanie tożsamości
├── base/             \PYGZsh{} rzeczywiste bazy danych
├── global/           \PYGZsh{} tabele systemowe
├── log/              \PYGZsh{} logi serwera
└── pg\PYGZus{}wal/           \PYGZsh{} Write\PYGZhy{}Ahead Log (WAL)
\end{sphinxVerbatim}

\sphinxAtStartPar
\sphinxstylestrong{Sprawdzenie aktualnej lokalizacji:}

\begin{sphinxVerbatim}[commandchars=\\\{\}]
\PYG{k}{SHOW}\PYG{+w}{ }\PYG{n}{data\PYGZus{}directory}\PYG{p}{;}
\PYG{k}{SHOW}\PYG{+w}{ }\PYG{n}{config\PYGZus{}file}\PYG{p}{;}
\PYG{k}{SHOW}\PYG{+w}{ }\PYG{n}{hba\PYGZus{}file}\PYG{p}{;}
\end{sphinxVerbatim}


\subsection{Środowiska i profile}
\label{\detokenize{rozdzial_2/rozdzial_22/index:srodowiska-i-profile}}
\sphinxAtStartPar
Przełączanie środowisk PostgreSQL (dev/test/prod) realizuj przez:
\begin{itemize}
\item {} 
\sphinxAtStartPar
zmienne środowiskowe (np. \sphinxcode{\sphinxupquote{PGHOST}}, \sphinxcode{\sphinxupquote{PGPORT}}, \sphinxcode{\sphinxupquote{PGDATABASE}}, \sphinxcode{\sphinxupquote{PGUSER}}, \sphinxcode{\sphinxupquote{PGPASSWORD}}),

\item {} 
\sphinxAtStartPar
profile konfiguracyjne,

\item {} 
\sphinxAtStartPar
osobne pliki \sphinxcode{\sphinxupquote{.env}} dla każdego środowiska.

\end{itemize}


\subsection{Automatyzacja}
\label{\detokenize{rozdzial_2/rozdzial_22/index:automatyzacja}}
\sphinxAtStartPar
Zalecenia dla PostgreSQL:
\begin{itemize}
\item {} 
\sphinxAtStartPar
walidacja konfiguracji przy starcie za pomocą \sphinxcode{\sphinxupquote{pg\_isready}},

\item {} 
\sphinxAtStartPar
skrypty sprawdzające kompletność zmiennych środowiskowych,

\item {} 
\sphinxAtStartPar
integracja z CI/CD (np. GitHub Actions, GitLab CI).

\end{itemize}

\sphinxAtStartPar
Dokumentację generuj automatycznie (Sphinx).


\subsection{Przykłady}
\label{\detokenize{rozdzial_2/rozdzial_22/index:przyklady}}
\sphinxAtStartPar
\sphinxstylestrong{PostgreSQL \textendash{} zmienne środowiskowe (.env)}

\begin{sphinxVerbatim}[commandchars=\\\{\}]
PGHOST=localhost
PGPORT=5432
PGDATABASE=aplikacja
PGUSER=app\PYGZus{}user
PGPASSWORD=\PYGZdl{}\PYGZob{}DB\PYGZus{}PASSWORD\PYGZcb{}
\end{sphinxVerbatim}

\sphinxAtStartPar
\sphinxstylestrong{PostgreSQL \textendash{} URL połączenia}

\begin{sphinxVerbatim}[commandchars=\\\{\}]
DATABASE\PYGZus{}URL=postgresql://\PYGZdl{}\PYGZob{}PGUSER\PYGZcb{}:\PYGZdl{}\PYGZob{}PGPASSWORD\PYGZcb{}@\PYGZdl{}\PYGZob{}PGHOST\PYGZcb{}:\PYGZdl{}\PYGZob{}PGPORT\PYGZcb{}/\PYGZdl{}\PYGZob{}PGDATABASE\PYGZcb{}
\end{sphinxVerbatim}

\sphinxAtStartPar
\sphinxstylestrong{PostgreSQL \textendash{} sprawdzenie połączenia}

\begin{sphinxVerbatim}[commandchars=\\\{\}]
pg\PYGZus{}isready \PYGZhy{}h \PYGZdl{}\PYGZob{}PGHOST\PYGZcb{} \PYGZhy{}p \PYGZdl{}\PYGZob{}PGPORT\PYGZcb{} \PYGZhy{}U \PYGZdl{}\PYGZob{}PGUSER\PYGZcb{} \PYGZhy{}d \PYGZdl{}\PYGZob{}PGDATABASE\PYGZcb{}
\end{sphinxVerbatim}


\subsection{Najlepsze praktyki}
\label{\detokenize{rozdzial_2/rozdzial_22/index:najlepsze-praktyki}}\begin{enumerate}
\sphinxsetlistlabels{\arabic}{enumi}{enumii}{}{.}%
\item {} 
\sphinxAtStartPar
Nie przechowuj sekretów w repozytorium.

\item {} 
\sphinxAtStartPar
Stosuj \sphinxcode{\sphinxupquote{.env.example}} zamiast rzeczywistych danych.

\item {} 
\sphinxAtStartPar
Używaj standardowych zmiennych środowiskowych PostgreSQL (\sphinxcode{\sphinxupquote{PG*}}).

\item {} 
\sphinxAtStartPar
Waliduj konfigurację przed uruchomieniem (\sphinxcode{\sphinxupquote{pg\_isready}}).

\item {} 
\sphinxAtStartPar
W środowiskach produkcyjnych używaj połączeń SSL/TLS (\sphinxcode{\sphinxupquote{sslmode=require}}).

\end{enumerate}


\subsection{Podsumowanie}
\label{\detokenize{rozdzial_2/rozdzial_22/index:podsumowanie}}
\sphinxAtStartPar
Poprawna konfiguracja połączenia z PostgreSQL zwiększa bezpieczeństwo i przenośność aplikacji między środowiskami. Standardowe zmienne środowiskowe oraz narzędzia takie jak \sphinxcode{\sphinxupquote{pg\_isready}} ułatwiają automatyzację i walidację.


\subsection{Szczegółowe omówienie plików konfiguracyjnych}
\label{\detokenize{rozdzial_2/rozdzial_22/index:szczegolowe-omowienie-plikow-konfiguracyjnych}}
\sphinxAtStartPar
Po omówieniu podstawowych wymagań środowiskowych można przejść do szczegółowej konfiguracji PostgreSQL.

\sphinxAtStartPar
W kolejnych sekcjach przedstawiono najważniejsze parametry konfiguracyjne dostępne w plikach:
\begin{itemize}
\item {} 
\sphinxAtStartPar
\sphinxcode{\sphinxupquote{postgresql.conf}} \textendash{} konfiguracja działania serwera,

\item {} 
\sphinxAtStartPar
\sphinxcode{\sphinxupquote{pg\_hba.conf}} \textendash{} kontrola dostępu do bazy danych,

\item {} 
\sphinxAtStartPar
\sphinxcode{\sphinxupquote{pg\_ident.conf}} \textendash{} mapowanie użytkowników systemowych na role PostgreSQL.

\end{itemize}


\subsection{postgresql.conf}
\label{\detokenize{rozdzial_2/rozdzial_22/index:postgresql-conf}}
\sphinxAtStartPar
Plik \sphinxcode{\sphinxupquote{postgresql.conf}} jest głównym plikiem konfiguracyjnym klastra PostgreSQL.
Zawiera ustawienia wpływające na działanie serwera, wydajność, bezpieczeństwo, logowanie oraz replikację.

\sphinxAtStartPar
Lokalizację aktywnego pliku można sprawdzić poleceniem:

\begin{sphinxVerbatim}[commandchars=\\\{\}]
\PYG{k}{SHOW}\PYG{+w}{ }\PYG{n}{config\PYGZus{}file}\PYG{p}{;}
\end{sphinxVerbatim}


\subsubsection{Najważniejsze parametry}
\label{\detokenize{rozdzial_2/rozdzial_22/index:najwazniejsze-parametry}}

\paragraph{Połączenia sieciowe}
\label{\detokenize{rozdzial_2/rozdzial_22/index:polaczenia-sieciowe}}\begin{description}
\sphinxlineitem{\sphinxcode{\sphinxupquote{listen\_addresses}}}
\sphinxAtStartPar
Określa adresy IP, na których PostgreSQL nasłuchuje połączeń.

\sphinxAtStartPar
Najczęstsze wartości:
\begin{itemize}
\item {} 
\sphinxAtStartPar
\sphinxcode{\sphinxupquote{localhost}} \textendash{} tylko połączenia lokalne,

\item {} 
\sphinxAtStartPar
\sphinxcode{\sphinxupquote{*}} \textendash{} wszystkie interfejsy sieciowe,

\item {} 
\sphinxAtStartPar
lista adresów, np. \sphinxcode{\sphinxupquote{\textquotesingle{}192.168.1.10,localhost\textquotesingle{}}}.

\end{itemize}

\sphinxlineitem{\sphinxcode{\sphinxupquote{port}}}
\sphinxAtStartPar
Port TCP używany przez PostgreSQL.

\sphinxAtStartPar
Domyślna wartość:

\begin{sphinxVerbatim}[commandchars=\\\{\}]
\PYG{n+na}{port}\PYG{+w}{ }\PYG{o}{=}\PYG{+w}{ }\PYG{l+s}{5432}
\end{sphinxVerbatim}

\sphinxlineitem{\sphinxcode{\sphinxupquote{max\_connections}}}
\sphinxAtStartPar
Maksymalna liczba jednoczesnych połączeń z bazą danych.

\sphinxAtStartPar
Zbyt wysoka wartość może zwiększać zużycie pamięci.

\end{description}


\paragraph{Pamięć i wydajność}
\label{\detokenize{rozdzial_2/rozdzial_22/index:pamiec-i-wydajnosc}}\begin{description}
\sphinxlineitem{\sphinxcode{\sphinxupquote{shared\_buffers}}}
\sphinxAtStartPar
Ilość pamięci RAM używanej przez PostgreSQL do buforowania danych.

\sphinxAtStartPar
Typowe wartości:
\begin{itemize}
\item {} 
\sphinxAtStartPar
128MB \textendash{} środowiska testowe,

\item {} 
\sphinxAtStartPar
25\% pamięci RAM \textendash{} środowiska produkcyjne.

\end{itemize}

\sphinxlineitem{\sphinxcode{\sphinxupquote{work\_mem}}}
\sphinxAtStartPar
Ilość pamięci wykorzystywanej przez operacje sortowania i zapytania.

\sphinxAtStartPar
Parametr jest przydzielany dla pojedynczej operacji.

\sphinxlineitem{\sphinxcode{\sphinxupquote{maintenance\_work\_mem}}}
\sphinxAtStartPar
Pamięć używana podczas operacji administracyjnych, np.:
\begin{itemize}
\item {} 
\sphinxAtStartPar
\sphinxcode{\sphinxupquote{VACUUM}},

\item {} 
\sphinxAtStartPar
\sphinxcode{\sphinxupquote{CREATE INDEX}},

\item {} 
\sphinxAtStartPar
\sphinxcode{\sphinxupquote{ALTER TABLE}}.

\end{itemize}

\sphinxlineitem{\sphinxcode{\sphinxupquote{effective\_cache\_size}}}
\sphinxAtStartPar
Szacowany rozmiar pamięci podręcznej dostępnej dla PostgreSQL.

\sphinxAtStartPar
Parametr pomaga optymalizatorowi zapytań wybierać odpowiednie plany wykonania.

\end{description}


\paragraph{Logowanie}
\label{\detokenize{rozdzial_2/rozdzial_22/index:logowanie}}\begin{description}
\sphinxlineitem{\sphinxcode{\sphinxupquote{logging\_collector}}}
\sphinxAtStartPar
Włącza zapisywanie logów do plików.

\sphinxAtStartPar
Dostępne wartości:
\begin{itemize}
\item {} 
\sphinxAtStartPar
\sphinxcode{\sphinxupquote{on}},

\item {} 
\sphinxAtStartPar
\sphinxcode{\sphinxupquote{off}}.

\end{itemize}

\sphinxlineitem{\sphinxcode{\sphinxupquote{log\_directory}}}
\sphinxAtStartPar
Katalog przechowywania logów.

\sphinxlineitem{\sphinxcode{\sphinxupquote{log\_filename}}}
\sphinxAtStartPar
Wzorzec nazw plików logów.

\sphinxAtStartPar
Przykład:

\begin{sphinxVerbatim}[commandchars=\\\{\}]
\PYG{n+na}{log\PYGZus{}filename}\PYG{+w}{ }\PYG{o}{=}\PYG{+w}{ }\PYG{l+s}{\PYGZsq{}postgresql\PYGZhy{}\PYGZpc{}Y\PYGZhy{}\PYGZpc{}m\PYGZhy{}\PYGZpc{}d.log\PYGZsq{}}
\end{sphinxVerbatim}

\sphinxlineitem{\sphinxcode{\sphinxupquote{log\_min\_duration\_statement}}}
\sphinxAtStartPar
Określa minimalny czas wykonania zapytania (w ms), po którym zostanie ono zapisane w logach.

\sphinxAtStartPar
Przykład:

\begin{sphinxVerbatim}[commandchars=\\\{\}]
\PYG{n+na}{log\PYGZus{}min\PYGZus{}duration\PYGZus{}statement}\PYG{+w}{ }\PYG{o}{=}\PYG{+w}{ }\PYG{l+s}{1000}
\end{sphinxVerbatim}

\sphinxAtStartPar
Wartość \sphinxcode{\sphinxupquote{1000}} oznacza logowanie zapytań trwających dłużej niż 1 sekundę.

\end{description}


\paragraph{Bezpieczeństwo}
\label{\detokenize{rozdzial_2/rozdzial_22/index:bezpieczenstwo}}\begin{description}
\sphinxlineitem{\sphinxcode{\sphinxupquote{password\_encryption}}}
\sphinxAtStartPar
Algorytm szyfrowania haseł użytkowników.

\sphinxAtStartPar
Zalecana wartość:

\begin{sphinxVerbatim}[commandchars=\\\{\}]
\PYG{n+na}{password\PYGZus{}encryption}\PYG{+w}{ }\PYG{o}{=}\PYG{+w}{ }\PYG{l+s}{scram\PYGZhy{}sha\PYGZhy{}256}
\end{sphinxVerbatim}

\sphinxlineitem{\sphinxcode{\sphinxupquote{ssl}}}
\sphinxAtStartPar
Włącza obsługę połączeń SSL/TLS.

\sphinxAtStartPar
Dostępne wartości:
\begin{itemize}
\item {} 
\sphinxAtStartPar
\sphinxcode{\sphinxupquote{on}},

\item {} 
\sphinxAtStartPar
\sphinxcode{\sphinxupquote{off}}.

\end{itemize}

\sphinxlineitem{\sphinxcode{\sphinxupquote{ssl\_cert\_file}} / \sphinxcode{\sphinxupquote{ssl\_key\_file}}}
\sphinxAtStartPar
Ścieżki do certyfikatu i klucza prywatnego SSL.

\end{description}


\paragraph{Replikacja i WAL}
\label{\detokenize{rozdzial_2/rozdzial_22/index:replikacja-i-wal}}\begin{description}
\sphinxlineitem{\sphinxcode{\sphinxupquote{wal\_level}}}
\sphinxAtStartPar
Określa poziom szczegółowości Write\sphinxhyphen{}Ahead Log (WAL).

\sphinxAtStartPar
Dostępne wartości:
\begin{itemize}
\item {} 
\sphinxAtStartPar
\sphinxcode{\sphinxupquote{minimal}},

\item {} 
\sphinxAtStartPar
\sphinxcode{\sphinxupquote{replica}},

\item {} 
\sphinxAtStartPar
\sphinxcode{\sphinxupquote{logical}}.

\end{itemize}

\sphinxlineitem{\sphinxcode{\sphinxupquote{max\_wal\_size}}}
\sphinxAtStartPar
Maksymalny rozmiar plików WAL przed wykonaniem checkpointu.

\sphinxlineitem{\sphinxcode{\sphinxupquote{archive\_mode}}}
\sphinxAtStartPar
Włącza archiwizację WAL.

\sphinxlineitem{\sphinxcode{\sphinxupquote{archive\_command}}}
\sphinxAtStartPar
Polecenie wykonywane podczas archiwizacji plików WAL.

\end{description}


\subsubsection{Przykładowa konfiguracja}
\label{\detokenize{rozdzial_2/rozdzial_22/index:przykladowa-konfiguracja}}
\begin{sphinxVerbatim}[commandchars=\\\{\}]
\PYG{n+na}{listen\PYGZus{}addresses}\PYG{+w}{ }\PYG{o}{=}\PYG{+w}{ }\PYG{l+s}{\PYGZsq{}*\PYGZsq{}}
\PYG{n+na}{port}\PYG{+w}{ }\PYG{o}{=}\PYG{+w}{ }\PYG{l+s}{5432}
\PYG{n+na}{max\PYGZus{}connections}\PYG{+w}{ }\PYG{o}{=}\PYG{+w}{ }\PYG{l+s}{200}

\PYG{n+na}{shared\PYGZus{}buffers}\PYG{+w}{ }\PYG{o}{=}\PYG{+w}{ }\PYG{l+s}{1GB}
\PYG{n+na}{work\PYGZus{}mem}\PYG{+w}{ }\PYG{o}{=}\PYG{+w}{ }\PYG{l+s}{16MB}
\PYG{n+na}{maintenance\PYGZus{}work\PYGZus{}mem}\PYG{+w}{ }\PYG{o}{=}\PYG{+w}{ }\PYG{l+s}{256MB}

\PYG{n+na}{logging\PYGZus{}collector}\PYG{+w}{ }\PYG{o}{=}\PYG{+w}{ }\PYG{l+s}{on}
\PYG{n+na}{log\PYGZus{}directory}\PYG{+w}{ }\PYG{o}{=}\PYG{+w}{ }\PYG{l+s}{\PYGZsq{}log\PYGZsq{}}

\PYG{n+na}{ssl}\PYG{+w}{ }\PYG{o}{=}\PYG{+w}{ }\PYG{l+s}{on}
\PYG{n+na}{password\PYGZus{}encryption}\PYG{+w}{ }\PYG{o}{=}\PYG{+w}{ }\PYG{l+s}{scram\PYGZhy{}sha\PYGZhy{}256}

\PYG{n+na}{wal\PYGZus{}level}\PYG{+w}{ }\PYG{o}{=}\PYG{+w}{ }\PYG{l+s}{replica}
\PYG{n+na}{max\PYGZus{}wal\PYGZus{}size}\PYG{+w}{ }\PYG{o}{=}\PYG{+w}{ }\PYG{l+s}{2GB}
\end{sphinxVerbatim}


\subsubsection{Weryfikacja konfiguracji}
\label{\detokenize{rozdzial_2/rozdzial_22/index:weryfikacja-konfiguracji}}
\sphinxAtStartPar
Po zmianie parametrów konfigurację można przeładować bez restartu serwera:

\begin{sphinxVerbatim}[commandchars=\\\{\}]
SELECT\PYG{+w}{ }pg\PYGZus{}reload\PYGZus{}conf\PYG{o}{(}\PYG{o}{)}\PYG{p}{;}
\end{sphinxVerbatim}

\sphinxAtStartPar
Niektóre parametry wymagają pełnego restartu PostgreSQL.

\sphinxAtStartPar
Aktualne wartości parametrów można sprawdzić poleceniem:

\begin{sphinxVerbatim}[commandchars=\\\{\}]
\PYG{k}{SHOW}\PYG{+w}{ }\PYG{n}{shared\PYGZus{}buffers}\PYG{p}{;}
\PYG{k}{SHOW}\PYG{+w}{ }\PYG{n}{wal\PYGZus{}level}\PYG{p}{;}
\PYG{k}{SHOW}\PYG{+w}{ }\PYG{n}{max\PYGZus{}connections}\PYG{p}{;}
\end{sphinxVerbatim}


\subsection{pg\_hba.conf}
\label{\detokenize{rozdzial_2/rozdzial_22/index:pg-hba-conf}}
\sphinxAtStartPar
Plik \sphinxcode{\sphinxupquote{pg\_hba.conf}} (Host\sphinxhyphen{}Based Authentication) odpowiada za kontrolę dostępu do serwera PostgreSQL.

\sphinxAtStartPar
Definiuje:
\begin{itemize}
\item {} 
\sphinxAtStartPar
kto może połączyć się z bazą danych,

\item {} 
\sphinxAtStartPar
z jakiego adresu IP,

\item {} 
\sphinxAtStartPar
do jakiej bazy danych,

\item {} 
\sphinxAtStartPar
przy użyciu jakiej metody uwierzytelniania.

\end{itemize}

\sphinxAtStartPar
Lokalizację aktywnego pliku można sprawdzić poleceniem:

\begin{sphinxVerbatim}[commandchars=\\\{\}]
\PYG{k}{SHOW}\PYG{+w}{ }\PYG{n}{hba\PYGZus{}file}\PYG{p}{;}
\end{sphinxVerbatim}


\subsubsection{Składnia wpisu}
\label{\detokenize{rozdzial_2/rozdzial_22/index:skladnia-wpisu}}
\sphinxAtStartPar
Każdy wpis w \sphinxcode{\sphinxupquote{pg\_hba.conf}} ma następującą strukturę:

\begin{sphinxVerbatim}[commandchars=\\\{\}]
\PYG{n}{TYPE}    \PYG{n}{DATABASE}    \PYG{n}{USER}    \PYG{n}{ADDRESS}    \PYG{n}{METHOD}
\end{sphinxVerbatim}

\sphinxAtStartPar
Znaczenie pól:
\begin{description}
\sphinxlineitem{\sphinxcode{\sphinxupquote{TYPE}}}
\sphinxAtStartPar
Typ połączenia.

\sphinxAtStartPar
Dostępne wartości:
\begin{itemize}
\item {} 
\sphinxAtStartPar
\sphinxcode{\sphinxupquote{local}} \textendash{} połączenia lokalne (Unix socket),

\item {} 
\sphinxAtStartPar
\sphinxcode{\sphinxupquote{host}} \textendash{} połączenia TCP/IP,

\item {} 
\sphinxAtStartPar
\sphinxcode{\sphinxupquote{hostssl}} \textendash{} połączenia TCP/IP z SSL,

\item {} 
\sphinxAtStartPar
\sphinxcode{\sphinxupquote{hostnossl}} \textendash{} połączenia TCP/IP bez SSL.

\end{itemize}

\sphinxlineitem{\sphinxcode{\sphinxupquote{DATABASE}}}
\sphinxAtStartPar
Nazwa bazy danych, do której użytkownik może się połączyć.

\sphinxAtStartPar
Możliwe wartości:
\begin{itemize}
\item {} 
\sphinxAtStartPar
konkretna nazwa bazy,

\item {} 
\sphinxAtStartPar
\sphinxcode{\sphinxupquote{all}} \textendash{} wszystkie bazy,

\item {} 
\sphinxAtStartPar
\sphinxcode{\sphinxupquote{sameuser}} \textendash{} baza o nazwie zgodnej z użytkownikiem.

\end{itemize}

\sphinxlineitem{\sphinxcode{\sphinxupquote{USER}}}
\sphinxAtStartPar
Użytkownik lub rola PostgreSQL.

\sphinxlineitem{\sphinxcode{\sphinxupquote{ADDRESS}}}
\sphinxAtStartPar
Adres IP lub zakres sieci w notacji CIDR.

\sphinxAtStartPar
Przykłady:
\begin{itemize}
\item {} 
\sphinxAtStartPar
\sphinxcode{\sphinxupquote{127.0.0.1/32}},

\item {} 
\sphinxAtStartPar
\sphinxcode{\sphinxupquote{192.168.1.0/24}},

\item {} 
\sphinxAtStartPar
\sphinxcode{\sphinxupquote{0.0.0.0/0}} \textendash{} wszystkie adresy.

\end{itemize}

\sphinxlineitem{\sphinxcode{\sphinxupquote{METHOD}}}
\sphinxAtStartPar
Metoda uwierzytelniania użytkownika.

\end{description}


\subsubsection{Najczęściej używane metody uwierzytelniania}
\label{\detokenize{rozdzial_2/rozdzial_22/index:najczesciej-uzywane-metody-uwierzytelniania}}\begin{description}
\sphinxlineitem{\sphinxcode{\sphinxupquote{trust}}}
\sphinxAtStartPar
Zezwala na połączenie bez uwierzytelniania.

\sphinxAtStartPar
Zalecane wyłącznie w środowiskach testowych.

\sphinxlineitem{\sphinxcode{\sphinxupquote{reject}}}
\sphinxAtStartPar
Odrzuca połączenie.

\sphinxlineitem{\sphinxcode{\sphinxupquote{md5}}}
\sphinxAtStartPar
Uwierzytelnianie przy użyciu hasła MD5.

\sphinxAtStartPar
Metoda starsza i mniej bezpieczna niż SCRAM.

\sphinxlineitem{\sphinxcode{\sphinxupquote{scram\sphinxhyphen{}sha\sphinxhyphen{}256}}}
\sphinxAtStartPar
Nowoczesna metoda uwierzytelniania oparta na SCRAM.

\sphinxAtStartPar
Zalecana dla nowych instalacji PostgreSQL.

\sphinxlineitem{\sphinxcode{\sphinxupquote{peer}}}
\sphinxAtStartPar
Uwierzytelnianie na podstawie użytkownika systemowego.

\sphinxAtStartPar
Najczęściej używane dla lokalnych połączeń Linux/Unix.

\sphinxlineitem{\sphinxcode{\sphinxupquote{ident}}}
\sphinxAtStartPar
Mapowanie użytkownika systemowego przy użyciu \sphinxcode{\sphinxupquote{pg\_ident.conf}}.

\end{description}


\subsubsection{Przykładowa konfiguracja}
\label{\detokenize{rozdzial_2/rozdzial_22/index:id1}}
\begin{sphinxVerbatim}[commandchars=\\\{\}]
\PYG{c+c1}{\PYGZsh{} Połączenia lokalne}
\PYG{n+na}{local   all             postgres                        peer}

\PYG{c+c1}{\PYGZsh{} Połączenia localhost IPv4}
\PYG{n+na}{host    all             all         127.0.0.1/32        scram\PYGZhy{}sha\PYGZhy{}256}

\PYG{c+c1}{\PYGZsh{} Połączenia localhost IPv6}
\PYG{n+na}{host    all             all}\PYG{+w}{         }\PYG{o}{:}\PYG{l+s}{:1/128             scram\PYGZhy{}sha\PYGZhy{}256}

\PYG{c+c1}{\PYGZsh{} Sieć lokalna}
\PYG{n+na}{host    aplikacja       app\PYGZus{}user    192.168.1.0/24      scram\PYGZhy{}sha\PYGZhy{}256}
\end{sphinxVerbatim}


\subsubsection{Najlepsze praktyki}
\label{\detokenize{rozdzial_2/rozdzial_22/index:id2}}\begin{enumerate}
\sphinxsetlistlabels{\arabic}{enumi}{enumii}{}{.}%
\item {} 
\sphinxAtStartPar
Używaj \sphinxcode{\sphinxupquote{scram\sphinxhyphen{}sha\sphinxhyphen{}256}} zamiast \sphinxcode{\sphinxupquote{md5}}.

\item {} 
\sphinxAtStartPar
Ograniczaj dostęp do konkretnych adresów IP.

\item {} 
\sphinxAtStartPar
Nie używaj \sphinxcode{\sphinxupquote{trust}} w środowiskach produkcyjnych.

\item {} 
\sphinxAtStartPar
Stosuj \sphinxcode{\sphinxupquote{hostssl}} dla połączeń zdalnych.

\item {} 
\sphinxAtStartPar
Po zmianie konfiguracji wykonuj reload konfiguracji:

\end{enumerate}

\begin{sphinxVerbatim}[commandchars=\\\{\}]
\PYG{k}{SELECT}\PYG{+w}{ }\PYG{n}{pg\PYGZus{}reload\PYGZus{}conf}\PYG{p}{(}\PYG{p}{)}\PYG{p}{;}
\end{sphinxVerbatim}


\subsection{pg\_ident.conf}
\label{\detokenize{rozdzial_2/rozdzial_22/index:pg-ident-conf}}
\sphinxAtStartPar
Plik \sphinxcode{\sphinxupquote{pg\_ident.conf}} umożliwia mapowanie użytkowników systemowych na role PostgreSQL.

\sphinxAtStartPar
Najczęściej używany jest razem z metodami uwierzytelniania:
\begin{itemize}
\item {} 
\sphinxAtStartPar
\sphinxcode{\sphinxupquote{peer}},

\item {} 
\sphinxAtStartPar
\sphinxcode{\sphinxupquote{ident}}.

\end{itemize}

\sphinxAtStartPar
Lokalizację aktywnego pliku można sprawdzić poleceniem:

\begin{sphinxVerbatim}[commandchars=\\\{\}]
\PYG{k}{SHOW}\PYG{+w}{ }\PYG{n}{ident\PYGZus{}file}\PYG{p}{;}
\end{sphinxVerbatim}


\subsubsection{Zastosowanie}
\label{\detokenize{rozdzial_2/rozdzial_22/index:zastosowanie}}
\sphinxAtStartPar
Mechanizm mapowania pozwala określić, który użytkownik systemowy może zostać powiązany z konkretną rolą PostgreSQL.

\sphinxAtStartPar
Jest to szczególnie przydatne:
\begin{itemize}
\item {} 
\sphinxAtStartPar
w środowiskach Linux/Unix,

\item {} 
\sphinxAtStartPar
przy automatyzacji administracyjnej,

\item {} 
\sphinxAtStartPar
dla lokalnych połączeń bez użycia haseł.

\end{itemize}


\subsubsection{Składnia wpisu}
\label{\detokenize{rozdzial_2/rozdzial_22/index:id3}}
\begin{sphinxVerbatim}[commandchars=\\\{\}]
\PYG{n}{MAPNAME}    \PYG{n}{SYSTEM}\PYG{o}{\PYGZhy{}}\PYG{n}{USERNAME}    \PYG{n}{PG}\PYG{o}{\PYGZhy{}}\PYG{n}{USERNAME}
\end{sphinxVerbatim}

\sphinxAtStartPar
Znaczenie pól:
\begin{description}
\sphinxlineitem{\sphinxcode{\sphinxupquote{MAPNAME}}}
\sphinxAtStartPar
Nazwa mapowania wykorzystywana w \sphinxcode{\sphinxupquote{pg\_hba.conf}}.

\sphinxlineitem{\sphinxcode{\sphinxupquote{SYSTEM\sphinxhyphen{}USERNAME}}}
\sphinxAtStartPar
Użytkownik systemu operacyjnego.

\sphinxlineitem{\sphinxcode{\sphinxupquote{PG\sphinxhyphen{}USERNAME}}}
\sphinxAtStartPar
Rola PostgreSQL, na którą użytkownik ma zostać zmapowany.

\end{description}


\subsubsection{Przykładowa konfiguracja}
\label{\detokenize{rozdzial_2/rozdzial_22/index:id4}}
\begin{sphinxVerbatim}[commandchars=\\\{\}]
\PYG{c+c1}{\PYGZsh{} MAPNAME    SYSTEM\PYGZhy{}USERNAME    PG\PYGZhy{}USERNAME}
\PYG{n+na}{admin\PYGZus{}map    postgres           postgres}
\PYG{n+na}{admin\PYGZus{}map    deploy             db\PYGZus{}admin}
\PYG{n+na}{app\PYGZus{}map      appuser            aplikacja}
\end{sphinxVerbatim}

\sphinxAtStartPar
Przykład użycia w \sphinxcode{\sphinxupquote{pg\_hba.conf}}:

\begin{sphinxVerbatim}[commandchars=\\\{\}]
\PYG{n+na}{local   all     all     peer map}\PYG{o}{=}\PYG{l+s}{admin\PYGZus{}map}
\end{sphinxVerbatim}

\sphinxAtStartPar
W powyższym przykładzie użytkownicy systemowi zdefiniowani w \sphinxcode{\sphinxupquote{admin\_map}} mogą logować się jako odpowiadające im role PostgreSQL.


\subsubsection{Najlepsze praktyki}
\label{\detokenize{rozdzial_2/rozdzial_22/index:id5}}\begin{enumerate}
\sphinxsetlistlabels{\arabic}{enumi}{enumii}{}{.}%
\item {} 
\sphinxAtStartPar
Ograniczaj mapowania wyłącznie do wymaganych użytkowników.

\item {} 
\sphinxAtStartPar
Nie mapuj użytkowników systemowych na role superuser bez uzasadnienia.

\item {} 
\sphinxAtStartPar
Stosuj osobne role administracyjne i aplikacyjne.

\item {} 
\sphinxAtStartPar
Dokumentuj wszystkie mapowania użytkowników.

\item {} 
\sphinxAtStartPar
Regularnie przeglądaj konfigurację dostępu.

\end{enumerate}

\sphinxstepscope


\section{Monitorowanie i diagnostyka}
\label{\detokenize{rozdzial_2/rozdzial_24/index:monitorowanie-i-diagnostyka}}\label{\detokenize{rozdzial_2/rozdzial_24/index::doc}}
\sphinxstepscope

\sphinxAtStartPar
Tego typu jakis tekst

\sphinxstepscope


\section{Partycjonowanie danych}
\label{\detokenize{rozdzial_2/rozdzial_26/index:partycjonowanie-danych}}\label{\detokenize{rozdzial_2/rozdzial_26/index::doc}}
\sphinxstepscope


\section{Bezpieczeństwo}
\label{\detokenize{rozdzial_2/rozdzial_27/index:bezpieczenstwo}}\label{\detokenize{rozdzial_2/rozdzial_27/index::doc}}
\sphinxstepscope


\section{Sprawozdanie z laboratorium documentation}
\label{\detokenize{rozdzial_2/rozdzial_28/source/index:sprawozdanie-z-laboratorium-documentation}}\label{\detokenize{rozdzial_2/rozdzial_28/source/index::doc}}
\sphinxAtStartPar
Add your content using \sphinxcode{\sphinxupquote{reStructuredText}} syntax. See the
\sphinxhref{https://www.sphinx-doc.org/en/master/usage/restructuredtext/index.html}{reStructuredText}
documentation for details.

\sphinxstepscope


\subsection{Wprowadzenie}
\label{\detokenize{rozdzial_2/rozdzial_28/source/rozdzial_1/index:wprowadzenie}}\label{\detokenize{rozdzial_2/rozdzial_28/source/rozdzial_1/index::doc}}
\sphinxstepscope


\subsection{Badania literaturowe}
\label{\detokenize{rozdzial_2/rozdzial_28/source/rozdzial_2/index:badania-literaturowe}}\label{\detokenize{rozdzial_2/rozdzial_28/source/rozdzial_2/index::doc}}
\sphinxstepscope


\subsection{Zadania}
\label{\detokenize{rozdzial_2/rozdzial_28/source/rozdzial_3/index:zadania}}\label{\detokenize{rozdzial_2/rozdzial_28/source/rozdzial_3/index::doc}}
\sphinxstepscope


\chapter{Rozdzial 3: Projektowanie Bazy Danych: System Zarządzania Biblioteką}
\label{\detokenize{rozdzial_3/index:rozdzial-3-projektowanie-bazy-danych-system-zarzadzania-biblioteka}}\label{\detokenize{rozdzial_3/index::doc}}\begin{quote}\begin{description}
\sphinxlineitem{Autorzy}\begin{enumerate}
\sphinxsetlistlabels{\arabic}{enumi}{enumii}{}{.}%
\item {} 
\sphinxAtStartPar
Paweł Łoćwin

\item {} 
\sphinxAtStartPar
Paweł Łosowski

\end{enumerate}

\end{description}\end{quote}


\section{Wybór zagadnienia, opis procesów i danych}
\label{\detokenize{rozdzial_3/index:wybor-zagadnienia-opis-procesow-i-danych}}

\subsection{Wybrane zagadnienie:}
\label{\detokenize{rozdzial_3/index:wybrane-zagadnienie}}
\sphinxAtStartPar
System zarządzania wypożyczeniami w bibliotece.
Projekt obejmuje obsługę bazy czytelników, katalogu zbiorów bibliotecznych (z uwzględnieniem autorów i kategorii literackich), a także pełen proces ewidencji wypożyczeń, zwrotów oraz historii aktywności czytelnika.


\subsection{Opis procesów i więzy integralności:}
\label{\detokenize{rozdzial_3/index:opis-procesow-i-wiezy-integralnosci}}
\sphinxAtStartPar
Głównym procesem w systemie jest obieg książki między biblioteką a czytelnikiem.
\begin{itemize}
\item {} 
\sphinxAtStartPar
\sphinxstylestrong{Rejestracja:} Czytelnik zapisuje się do biblioteki, podając swoje dane osobowe oraz adresowe. System automatycznie rejestruje datę zapisu.

\item {} 
\sphinxAtStartPar
\sphinxstylestrong{Katalogowanie:} Książki są kategoryzowane (np. Fantastyka, Kryminał, Literatura faktu) i przypisane do konkretnych autorów. Każdy fizyczny egzemplarz książki posiada swój unikalny identyfikator.

\item {} 
\sphinxAtStartPar
\sphinxstylestrong{Wypożyczenia i Zwroty:} Proces wypożyczenia łączy czytelnika z konkretnym egzemplarzem książki w czasie. Rejestrowana jest data wypożyczenia. System zakłada konieczność ewidencji daty zwrotu, by w przyszłości móc naliczać ewentualne kary.

\item {} 
\sphinxAtStartPar
\sphinxstylestrong{Statystyki (Więzy integralności):} Pola takie jak “aktualne wypożyczenia” i “suma wypożyczeń” z pierwotnego prototypu są wartościami wyliczalnymi (pochodnymi) na podstawie historii transakcji, co w docelowym modelu gwarantuje spójność danych.

\end{itemize}


\subsection{Wykaz gromadzonych danych:}
\label{\detokenize{rozdzial_3/index:wykaz-gromadzonych-danych}}\begin{itemize}
\item {} 
\sphinxAtStartPar
\sphinxstylestrong{Dane czytelnika:} Imię, Nazwisko, Ulica, Kod pocztowy, Miasto, Data zapisu.

\item {} 
\sphinxAtStartPar
\sphinxstylestrong{Dane książki:} Tytuł, Rok wydania, Dostępność.

\item {} 
\sphinxAtStartPar
\sphinxstylestrong{Dane autora:} Imię, Nazwisko.

\item {} 
\sphinxAtStartPar
\sphinxstylestrong{Dane kategorii:} Nazwa gatunku literackiego.

\item {} 
\sphinxAtStartPar
\sphinxstylestrong{Dane transakcyjne:} Data wypożyczenia, Data zwrotu (rzeczywista).

\end{itemize}


\section{Prototyp CSV}
\label{\detokenize{rozdzial_3/index:prototyp-csv}}
\sphinxAtStartPar
Aby zweryfikować kompletność przetwarzanych informacji, przygotowano “płaską” (nieznormalizowaną) reprezentację danych dla operacji wypożyczenia, bazującą na pierwotnych wytycznych.

\begin{sphinxVerbatim}[commandchars=\\\{\}]
\PYG{n}{Imie}\PYG{p}{,}\PYG{n}{Nazwisko}\PYG{p}{,}\PYG{n}{Ulica}\PYG{p}{,}\PYG{n}{Kod\PYGZus{}Pocztowy}\PYG{p}{,}\PYG{n}{Miasto}\PYG{p}{,}\PYG{n}{Data\PYGZus{}Zapisu}\PYG{p}{,}\PYG{n}{Tytul\PYGZus{}Ksiazki}\PYG{p}{,}\PYG{n}{Imie\PYGZus{}Autora}\PYG{p}{,}\PYG{n}{Nazwisko\PYGZus{}Autora}\PYG{p}{,}\PYG{n}{Kategoria}\PYG{p}{,}\PYG{n}{Data\PYGZus{}Wypozyczenia}\PYG{p}{,}\PYG{n}{Data\PYGZus{}Zwrotu}
\PYG{n}{Jan}\PYG{p}{,}\PYG{n}{Kowalski}\PYG{p}{,}\PYG{n}{Kwiatowa} \PYG{l+m+mi}{1}\PYG{p}{,}\PYG{l+m+mi}{00}\PYG{o}{\PYGZhy{}}\PYG{l+m+mi}{001}\PYG{p}{,}\PYG{n}{Warszawa}\PYG{p}{,}\PYG{l+m+mi}{2024}\PYG{o}{\PYGZhy{}}\PYG{l+m+mi}{01}\PYG{o}{\PYGZhy{}}\PYG{l+m+mi}{15}\PYG{p}{,}\PYG{n}{Wiedźmin}\PYG{p}{,}\PYG{n}{Andrzej}\PYG{p}{,}\PYG{n}{Sapkowski}\PYG{p}{,}\PYG{n}{Fantastyka}\PYG{p}{,}\PYG{l+m+mi}{2024}\PYG{o}{\PYGZhy{}}\PYG{l+m+mi}{03}\PYG{o}{\PYGZhy{}}\PYG{l+m+mi}{01}\PYG{p}{,}\PYG{l+m+mi}{2024}\PYG{o}{\PYGZhy{}}\PYG{l+m+mi}{03}\PYG{o}{\PYGZhy{}}\PYG{l+m+mi}{15}
\PYG{n}{Jan}\PYG{p}{,}\PYG{n}{Kowalski}\PYG{p}{,}\PYG{n}{Kwiatowa} \PYG{l+m+mi}{1}\PYG{p}{,}\PYG{l+m+mi}{00}\PYG{o}{\PYGZhy{}}\PYG{l+m+mi}{001}\PYG{p}{,}\PYG{n}{Warszawa}\PYG{p}{,}\PYG{l+m+mi}{2024}\PYG{o}{\PYGZhy{}}\PYG{l+m+mi}{01}\PYG{o}{\PYGZhy{}}\PYG{l+m+mi}{15}\PYG{p}{,}\PYG{n}{Pan} \PYG{n}{Tadeusz}\PYG{p}{,}\PYG{n}{Adam}\PYG{p}{,}\PYG{n}{Mickiewicz}\PYG{p}{,}\PYG{n}{Poezja}\PYG{p}{,}\PYG{l+m+mi}{2024}\PYG{o}{\PYGZhy{}}\PYG{l+m+mi}{03}\PYG{o}{\PYGZhy{}}\PYG{l+m+mi}{10}\PYG{p}{,}
\PYG{n}{Anna}\PYG{p}{,}\PYG{n}{Nowak}\PYG{p}{,}\PYG{n}{Polna} \PYG{l+m+mi}{2}\PYG{p}{,}\PYG{l+m+mi}{31}\PYG{o}{\PYGZhy{}}\PYG{l+m+mi}{444}\PYG{p}{,}\PYG{n}{Kraków}\PYG{p}{,}\PYG{l+m+mi}{2023}\PYG{o}{\PYGZhy{}}\PYG{l+m+mi}{11}\PYG{o}{\PYGZhy{}}\PYG{l+m+mi}{05}\PYG{p}{,}\PYG{n}{Lśnienie}\PYG{p}{,}\PYG{n}{Stephen}\PYG{p}{,}\PYG{n}{King}\PYG{p}{,}\PYG{n}{Horror}\PYG{p}{,}\PYG{l+m+mi}{2024}\PYG{o}{\PYGZhy{}}\PYG{l+m+mi}{02}\PYG{o}{\PYGZhy{}}\PYG{l+m+mi}{20}\PYG{p}{,}\PYG{l+m+mi}{2024}\PYG{o}{\PYGZhy{}}\PYG{l+m+mi}{03}\PYG{o}{\PYGZhy{}}\PYG{l+m+mi}{05}
\end{sphinxVerbatim}


\section{Model Konceptualny (Pojęciowy)}
\label{\detokenize{rozdzial_3/index:model-konceptualny-pojeciowy}}
\sphinxAtStartPar
Na podstawie analizy procesów i zebranych danych opracowano model pojęciowy, identyfikując obiekty, ich cechy oraz powiązania.


\subsection{Zidentyfikowane encje}
\label{\detokenize{rozdzial_3/index:zidentyfikowane-encje}}\begin{itemize}
\item {} 
\sphinxAtStartPar
\sphinxstylestrong{Czytelnik} \sphinxhyphen{} osoba zapisana do biblioteki.

\item {} 
\sphinxAtStartPar
\sphinxstylestrong{Autor} \sphinxhyphen{} twórca dzieła literackiego.

\item {} 
\sphinxAtStartPar
\sphinxstylestrong{Książka} \sphinxhyphen{} oferowana pozycja z inwentarza biblioteki.

\item {} 
\sphinxAtStartPar
\sphinxstylestrong{Kategoria} \sphinxhyphen{} sekcja grupująca gatunkowo księgozbiór.

\item {} 
\sphinxAtStartPar
\sphinxstylestrong{Wypożyczenie} \sphinxhyphen{} zdarzenie potwierdzające fakt wydania książki czytelnikowi.

\end{itemize}


\subsection{Zdefiniowane atrybuty/własności}
\label{\detokenize{rozdzial_3/index:zdefiniowane-atrybuty-wlasnosci}}\begin{itemize}
\item {} 
\sphinxAtStartPar
\sphinxstylestrong{Czytelnik:} Imię, Nazwisko, Ulica, Kod pocztowy, Miasto, Data zapisu.

\item {} 
\sphinxAtStartPar
\sphinxstylestrong{Autor:} Imię, Nazwisko.

\item {} 
\sphinxAtStartPar
\sphinxstylestrong{Książka:} Tytuł, Rok wydania.

\item {} 
\sphinxAtStartPar
\sphinxstylestrong{Kategoria:} Nazwa kategorii.

\item {} 
\sphinxAtStartPar
\sphinxstylestrong{Wypożyczenie:} Data wypożyczenia, Data zwrotu.

\end{itemize}


\subsection{Opis związków}
\label{\detokenize{rozdzial_3/index:opis-zwiazkow}}\begin{itemize}
\item {} 
\sphinxAtStartPar
\sphinxstylestrong{Autor \textendash{} Książka (1:N):} Jeden autor może napisać wiele książek. (Dla uproszczenia modelu zakładamy głównego autora).

\item {} 
\sphinxAtStartPar
\sphinxstylestrong{Kategoria \textendash{} Książka (1:N):} Kategoria zawiera wiele tytułów.

\item {} 
\sphinxAtStartPar
\sphinxstylestrong{Czytelnik \textendash{} Wypożyczenie (1:N):} Czytelnik może dokonać wielu wypożyczeń w historii.

\item {} 
\sphinxAtStartPar
\sphinxstylestrong{Książka \textendash{} Wypożyczenie (1:N):} Jeden egzemplarz książki może być w swojej historii wypożyczany wielokrotnie różnym czytelnikom.

\end{itemize}


\subsection{Identyfikacja encji słabych}
\label{\detokenize{rozdzial_3/index:identyfikacja-encji-slabych}}
\sphinxAtStartPar
Występuje relacja wiele\sphinxhyphen{}do\sphinxhyphen{}wielu (M:N) między Czytelnikiem a Książką (czytelnik czyta wiele książek, książka jest czytana przez wielu czytelników). Związek ten rozwiązywany jest przez encję asocjacyjną (słabą) \sphinxstylestrong{Wypożyczenie}.
* \sphinxstylestrong{Uzasadnienie:} Byt ten nie istnieje samodzielnie w oderwaniu od konkretnego Czytelnika i konkretnej Książki.


\subsection{Schemat w notacji Chena}
\label{\detokenize{rozdzial_3/index:schemat-w-notacji-chena}}
\noindent\sphinxincludegraphics{{model_konceptualny_biblioteka1}.png}


\section{Model logiczny i proces normalizacji}
\label{\detokenize{rozdzial_3/index:model-logiczny-i-proces-normalizacji}}

\subsection{Przebieg procesu normalizacji}
\label{\detokenize{rozdzial_3/index:przebieg-procesu-normalizacji}}
\sphinxAtStartPar
\sphinxstylestrong{Krok 1: Pierwsza Postać Normalna (1NF)}
Wiersze są unikalne, tworzymy sztuczne klucze główne (ID\_Wypozyczenia). Wartości w komórkach są atomowe. Pojawia się duża redundancja danych adresowych i książkowych.

\sphinxAtStartPar
\sphinxstylestrong{Krok 2: Druga Postać Normalna (2NF)}
Rozbijamy płaską strukturę względem częściowych zależności. Oddzielamy encje twarde od operacji.
Tworzymy bazowe tabele: \sphinxcode{\sphinxupquote{Czytelnicy}} oraz \sphinxcode{\sphinxupquote{Ksiazki}}. Powstaje tabela \sphinxcode{\sphinxupquote{Wypozyczenia}} przechowująca klucze obce ID\_Czytelnika oraz ID\_Ksiazki, a także daty transakcji.
\sphinxstyleemphasis{Uwaga:} Dynamiczne atrybuty takie jak “suma wypożyczeń” z pierwotnego zadania są usuwane ze schematu, ponieważ łamią zasady normalizacji \textendash{} można je obliczyć zapytaniem SQL typu \sphinxcode{\sphinxupquote{COUNT()}}.

\sphinxAtStartPar
\sphinxstylestrong{Krok 3: Trzecia Postać Normalna (3NF)}
Eliminacja zależności przechodnich. Z tabeli \sphinxcode{\sphinxupquote{Ksiazki}} wydzielamy powtarzające się nazwy gatunków do tabeli \sphinxcode{\sphinxupquote{Kategorie}} (klucz: ID\_Kategorii) oraz dane twórców do tabeli \sphinxcode{\sphinxupquote{Autorzy}} (klucz: ID\_Autora).


\subsection{Ostateczna struktura tabel (3NF)}
\label{\detokenize{rozdzial_3/index:ostateczna-struktura-tabel-3nf}}\begin{itemize}
\item {} 
\sphinxAtStartPar
\sphinxstylestrong{Czytelnicy:} ID\_Czytelnika (PK), Imie, Nazwisko, Ulica, Kod\_Pocztowy, Miasto, Data\_Zapisu

\item {} 
\sphinxAtStartPar
\sphinxstylestrong{Autorzy:} ID\_Autora (PK), Imie, Nazwisko

\item {} 
\sphinxAtStartPar
\sphinxstylestrong{Kategorie:} ID\_Kategorii (PK), Nazwa\_Kategorii

\item {} 
\sphinxAtStartPar
\sphinxstylestrong{Ksiazki:} ID\_Ksiazki (PK), ID\_Autora (FK), ID\_Kategorii (FK), Tytul, Rok\_Wydania

\item {} 
\sphinxAtStartPar
\sphinxstylestrong{Wypozyczenia:} ID\_Wypozyczenia (PK), ID\_Czytelnika (FK), ID\_Ksiazki (FK), Data\_Wypozyczenia, Data\_Zwrotu

\end{itemize}


\subsection{Diagram ERD (Model Logiczny)}
\label{\detokenize{rozdzial_3/index:diagram-erd-model-logiczny}}
\noindent\sphinxincludegraphics{{erd_logiczny_biblioteka1}.png}


\section{Model fizyczny bazy danych}
\label{\detokenize{rozdzial_3/index:model-fizyczny-bazy-danych}}
\sphinxAtStartPar
Różnice implementacyjne między modelami fizycznymi wynikają wprost z silników RDBMS \textendash{} ograniczeń typowania SQLite oraz bogatych możliwości deklaratywnych PostgreSQL.


\subsection{Model fizyczny dla środowiska SQLite}
\label{\detokenize{rozdzial_3/index:model-fizyczny-dla-srodowiska-sqlite}}
\sphinxAtStartPar
Z uwagi na okrojony zestaw typów, daty są mapowane jako TEXT.
\begin{itemize}
\item {} 
\sphinxAtStartPar
\sphinxstylestrong{Czytelnicy:} ID\_Czytelnika : INTEGER PRIMARY KEY, Imie : TEXT, Nazwisko : TEXT, Ulica : TEXT, Kod\_Pocztowy : TEXT, Miasto : TEXT, Data\_Zapisu : TEXT

\item {} 
\sphinxAtStartPar
\sphinxstylestrong{Autorzy:} ID\_Autora : INTEGER PRIMARY KEY, Imie : TEXT, Nazwisko : TEXT

\item {} 
\sphinxAtStartPar
\sphinxstylestrong{Kategorie:} ID\_Kategorii : INTEGER PRIMARY KEY, Nazwa\_Kategorii : TEXT

\item {} 
\sphinxAtStartPar
\sphinxstylestrong{Ksiazki:} ID\_Ksiazki : INTEGER PRIMARY KEY, ID\_Autora : INTEGER, ID\_Kategorii : INTEGER, Tytul : TEXT, Rok\_Wydania : INTEGER

\item {} 
\sphinxAtStartPar
\sphinxstylestrong{Wypozyczenia:} ID\_Wypozyczenia : INTEGER PRIMARY KEY, ID\_Czytelnika : INTEGER, ID\_Ksiazki : INTEGER, Data\_Wypozyczenia : TEXT, Data\_Zwrotu : TEXT

\end{itemize}

\noindent\sphinxincludegraphics{{fizyczny_sqlite_biblioteka1}.png}


\subsection{Model fizyczny dla środowiska PostgreSQL}
\label{\detokenize{rozdzial_3/index:model-fizyczny-dla-srodowiska-postgresql}}
\sphinxAtStartPar
PostgreSQL umożliwia zastosowanie precyzyjnych i natywnych typów, w tym rygorystycznych typów daty (DATE) oraz optymalizacji pamięciowej dla ciągów znaków (VARCHAR).
\begin{itemize}
\item {} 
\sphinxAtStartPar
\sphinxstylestrong{Czytelnicy:} ID\_Czytelnika : SERIAL PRIMARY KEY, Imie : VARCHAR(50), Nazwisko : VARCHAR(50), Ulica : VARCHAR(100), Kod\_Pocztowy : VARCHAR(6), Miasto : VARCHAR(50), Data\_Zapisu : DATE DEFAULT CURRENT\_DATE

\item {} 
\sphinxAtStartPar
\sphinxstylestrong{Autorzy:} ID\_Autora : SERIAL PRIMARY KEY, Imie : VARCHAR(50), Nazwisko : VARCHAR(50)

\item {} 
\sphinxAtStartPar
\sphinxstylestrong{Kategorie:} ID\_Kategorii : SERIAL PRIMARY KEY, Nazwa\_Kategorii : VARCHAR(50)

\item {} 
\sphinxAtStartPar
\sphinxstylestrong{Ksiazki:} ID\_Ksiazki : SERIAL PRIMARY KEY, ID\_Autora : INTEGER REFERENCES Autorzy, ID\_Kategorii : INTEGER REFERENCES Kategorie, Tytul : VARCHAR(150), Rok\_Wydania : SMALLINT

\item {} 
\sphinxAtStartPar
\sphinxstylestrong{Wypozyczenia:} ID\_Wypozyczenia : SERIAL PRIMARY KEY, ID\_Czytelnika : INTEGER REFERENCES Czytelnicy, ID\_Ksiazki : INTEGER REFERENCES Ksiazki, Data\_Wypozyczenia : DATE DEFAULT CURRENT\_DATE, Data\_Zwrotu : DATE

\end{itemize}

\noindent\sphinxincludegraphics{{fizyczny_postgres_biblioteka.drawio1}.png}

\sphinxstepscope


\chapter{Rozdział 4: Implementacja fizyczna i zasilanie bazy danych}
\label{\detokenize{rozdzial_4/index:rozdzial-4-implementacja-fizyczna-i-zasilanie-bazy-danych}}\label{\detokenize{rozdzial_4/index::doc}}\begin{quote}\begin{description}
\sphinxlineitem{Autorzy}\begin{enumerate}
\sphinxsetlistlabels{\arabic}{enumi}{enumii}{}{.}%
\item {} 
\sphinxAtStartPar
Paweł Łoćwin

\item {} 
\sphinxAtStartPar
Paweł Łosowski

\end{enumerate}

\end{description}\end{quote}


\section{Definicja fizyczna bazy danych (Zadanie 1)}
\label{\detokenize{rozdzial_4/index:definicja-fizyczna-bazy-danych-zadanie-1}}
\sphinxAtStartPar
Na podstawie opracowanych wcześniej modeli logicznych przygotowano skrypty DDL (Data Definition Language) dla dwóch docelowych silników relacyjnych baz danych: PostgreSQL oraz SQLite. Kody te posłużyły do utworzenia struktur w bazach dostępnych na serwerze laboratoryjnym oraz przez sieć Internet.


\subsection{Skrypt dla środowiska PostgreSQL}
\label{\detokenize{rozdzial_4/index:skrypt-dla-srodowiska-postgresql}}
\sphinxAtStartPar
Środowisko PostgreSQL pozwala na wykorzystanie zaawansowanych, natywnych typów danych oraz rygorystyczne egzekwowanie więzów integralności. W poniższym skrypcie kluczowym elementem jest użycie pseudo\sphinxhyphen{}typu \sphinxcode{\sphinxupquote{SERIAL}} dla automatycznej generacji kluczy głównych oraz typów o stałej lub ograniczonej długości (\sphinxcode{\sphinxupquote{VARCHAR}}, \sphinxcode{\sphinxupquote{SMALLINT}}, \sphinxcode{\sphinxupquote{DATE}}) dla optymalizacji przechowywania danych. Relacje między tabelami zapewniane są przez twarde klauzule \sphinxcode{\sphinxupquote{REFERENCES}}.

\begin{sphinxVerbatim}[commandchars=\\\{\}]
\PYG{c+c1}{\PYGZhy{}\PYGZhy{} Fragment kodu definicji kluczowych tabel (PostgreSQL)}
\PYG{k}{CREATE}\PYG{+w}{ }\PYG{k}{TABLE}\PYG{+w}{ }\PYG{n}{Ksiazki}\PYG{+w}{ }\PYG{p}{(}
\PYG{+w}{    }\PYG{n}{ID\PYGZus{}Ksiazki}\PYG{+w}{ }\PYG{n+nb}{SERIAL}\PYG{+w}{ }\PYG{k}{PRIMARY}\PYG{+w}{ }\PYG{k}{KEY}\PYG{p}{,}
\PYG{+w}{    }\PYG{n}{ID\PYGZus{}Autora}\PYG{+w}{ }\PYG{n+nb}{INTEGER}\PYG{+w}{ }\PYG{k}{REFERENCES}\PYG{+w}{ }\PYG{n}{Autorzy}\PYG{p}{(}\PYG{n}{ID\PYGZus{}Autora}\PYG{p}{)}\PYG{p}{,}
\PYG{+w}{    }\PYG{n}{ID\PYGZus{}Kategorii}\PYG{+w}{ }\PYG{n+nb}{INTEGER}\PYG{+w}{ }\PYG{k}{REFERENCES}\PYG{+w}{ }\PYG{n}{Kategorie}\PYG{p}{(}\PYG{n}{ID\PYGZus{}Kategorii}\PYG{p}{)}\PYG{p}{,}
\PYG{+w}{    }\PYG{n}{Tytul}\PYG{+w}{ }\PYG{n+nb}{VARCHAR}\PYG{p}{(}\PYG{l+m+mi}{150}\PYG{p}{)}\PYG{+w}{ }\PYG{k}{NOT}\PYG{+w}{ }\PYG{k}{NULL}\PYG{p}{,}
\PYG{+w}{    }\PYG{n}{Rok\PYGZus{}Wydania}\PYG{+w}{ }\PYG{n+nb}{SMALLINT}
\PYG{p}{)}\PYG{p}{;}

\PYG{k}{CREATE}\PYG{+w}{ }\PYG{k}{TABLE}\PYG{+w}{ }\PYG{n}{Wypozyczenia}\PYG{+w}{ }\PYG{p}{(}
\PYG{+w}{    }\PYG{n}{ID\PYGZus{}Wypozyczenia}\PYG{+w}{ }\PYG{n+nb}{SERIAL}\PYG{+w}{ }\PYG{k}{PRIMARY}\PYG{+w}{ }\PYG{k}{KEY}\PYG{p}{,}
\PYG{+w}{    }\PYG{n}{ID\PYGZus{}Czytelnika}\PYG{+w}{ }\PYG{n+nb}{INTEGER}\PYG{+w}{ }\PYG{k}{REFERENCES}\PYG{+w}{ }\PYG{n}{Czytelnicy}\PYG{p}{(}\PYG{n}{ID\PYGZus{}Czytelnika}\PYG{p}{)}\PYG{p}{,}
\PYG{+w}{    }\PYG{n}{ID\PYGZus{}Ksiazki}\PYG{+w}{ }\PYG{n+nb}{INTEGER}\PYG{+w}{ }\PYG{k}{REFERENCES}\PYG{+w}{ }\PYG{n}{Ksiazki}\PYG{p}{(}\PYG{n}{ID\PYGZus{}Ksiazki}\PYG{p}{)}\PYG{p}{,}
\PYG{+w}{    }\PYG{n}{Data\PYGZus{}Wypozyczenia}\PYG{+w}{ }\PYG{n+nb}{DATE}\PYG{+w}{ }\PYG{k}{DEFAULT}\PYG{+w}{ }\PYG{k}{CURRENT\PYGZus{}DATE}\PYG{p}{,}
\PYG{+w}{    }\PYG{n}{Data\PYGZus{}Zwrotu}\PYG{+w}{ }\PYG{n+nb}{DATE}
\PYG{p}{)}\PYG{p}{;}
\end{sphinxVerbatim}


\subsection{Skrypt dla środowiska SQLite}
\label{\detokenize{rozdzial_4/index:skrypt-dla-srodowiska-sqlite}}
\sphinxAtStartPar
W silniku SQLite struktura ulega pewnemu uproszczeniu z powodu mniejszej rygorystyczności typowania oraz bardzo ograniczonej liczby wbudowanych klas typów danych (np. brak odrębnego typu daty). Z tego względu zastosowano typ \sphinxcode{\sphinxupquote{TEXT}} dla dat i łańcuchów znaków. Ponadto w SQLite relacje kluczy obcych deklaruje się w osobnej klauzuli \sphinxcode{\sphinxupquote{FOREIGN KEY}} na końcu definicji tabeli, a automatyczna inkrementacja klucza realizowana jest atrybutem \sphinxcode{\sphinxupquote{AUTOINCREMENT}}.

\begin{sphinxVerbatim}[commandchars=\\\{\}]
\PYG{c+c1}{\PYGZhy{}\PYGZhy{} Fragment kodu definicji kluczowych tabel (SQLite)}
\PYG{k}{CREATE}\PYG{+w}{ }\PYG{k}{TABLE}\PYG{+w}{ }\PYG{n}{Wypozyczenia}\PYG{+w}{ }\PYG{p}{(}
\PYG{+w}{    }\PYG{n}{ID\PYGZus{}Wypozyczenia}\PYG{+w}{ }\PYG{n+nb}{INTEGER}\PYG{+w}{ }\PYG{k}{PRIMARY}\PYG{+w}{ }\PYG{k}{KEY}\PYG{+w}{ }\PYG{n}{AUTOINCREMENT}\PYG{p}{,}
\PYG{+w}{    }\PYG{n}{ID\PYGZus{}Czytelnika}\PYG{+w}{ }\PYG{n+nb}{INTEGER}\PYG{p}{,}
\PYG{+w}{    }\PYG{n}{ID\PYGZus{}Ksiazki}\PYG{+w}{ }\PYG{n+nb}{INTEGER}\PYG{p}{,}
\PYG{+w}{    }\PYG{n}{Data\PYGZus{}Wypozyczenia}\PYG{+w}{ }\PYG{n+nb}{TEXT}\PYG{p}{,}
\PYG{+w}{    }\PYG{n}{Data\PYGZus{}Zwrotu}\PYG{+w}{ }\PYG{n+nb}{TEXT}\PYG{p}{,}
\PYG{+w}{    }\PYG{k}{FOREIGN}\PYG{+w}{ }\PYG{k}{KEY}\PYG{p}{(}\PYG{n}{ID\PYGZus{}Czytelnika}\PYG{p}{)}\PYG{+w}{ }\PYG{k}{REFERENCES}\PYG{+w}{ }\PYG{n}{Czytelnicy}\PYG{p}{(}\PYG{n}{ID\PYGZus{}Czytelnika}\PYG{p}{)}\PYG{p}{,}
\PYG{+w}{    }\PYG{k}{FOREIGN}\PYG{+w}{ }\PYG{k}{KEY}\PYG{p}{(}\PYG{n}{ID\PYGZus{}Ksiazki}\PYG{p}{)}\PYG{+w}{ }\PYG{k}{REFERENCES}\PYG{+w}{ }\PYG{n}{Ksiazki}\PYG{p}{(}\PYG{n}{ID\PYGZus{}Ksiazki}\PYG{p}{)}
\PYG{p}{)}\PYG{p}{;}
\end{sphinxVerbatim}


\section{Mechanizm zasilania bazy danych (Zadanie 2)}
\label{\detokenize{rozdzial_4/index:mechanizm-zasilania-bazy-danych-zadanie-2}}
\sphinxAtStartPar
Posiadając gotową strukturę bazy, należało wyposażyć ją w dane demonstracyjne. Przygotowano zbiór danych testowych, które zostały poddane wsadowemu ładowaniu (batch import).


\subsection{Formatyzacja danych (Pliki CSV)}
\label{\detokenize{rozdzial_4/index:formatyzacja-danych-pliki-csv}}
\sphinxAtStartPar
Dane do importu przygotowano w niezależnym formacie płaskim CSV (Comma\sphinxhyphen{}Separated Values). Takie podejście pozwala na łatwą edycję danych poza systemem bazodanowym. Poniżej zaprezentowano wycinek pliku \sphinxcode{\sphinxupquote{dane\_kategorie.csv}}:

\begin{sphinxVerbatim}[commandchars=\\\{\}]
Nazwa\PYGZus{}Kategorii
Fantastyka
Kryminał
Literatura faktu
Poezja
Horror
\end{sphinxVerbatim}


\subsection{Implementacja mechanizmu importu}
\label{\detokenize{rozdzial_4/index:implementacja-mechanizmu-importu}}
\sphinxAtStartPar
Do zaimportowania powyższych danych wybrano mechanizm wprowadzania wielowartościowego oparty na technice \sphinxstylestrong{INSERT (multi\sphinxhyphen{}value/batch)}. Rozwiązanie zrealizowano z wykorzystaniem języka Python oraz dedykowanego sterownika \sphinxcode{\sphinxupquote{psycopg}}.

\sphinxAtStartPar
Wybór tego mechanizmu podyktowany jest kompromisem między uniwersalnością a wydajnością. Użycie metody \sphinxcode{\sphinxupquote{executemany()}} na obiekcie kursora bazy danych pozwala na szybkie wstawienie wielu rekordów w ramach pojedynczej transakcji sieciowej. Jest to rozwiązanie znacznie bardziej optymalne czasowo i obciążeniowo niż wykonywanie pojedynczych instrukcji \sphinxcode{\sphinxupquote{INSERT}} iteracyjnie w pętli. Opcjonalna klauzula \sphinxcode{\sphinxupquote{COPY}} byłaby szybsza, jednak rozwiązanie oparte o wsadowy \sphinxcode{\sphinxupquote{INSERT}} jest bardziej uniwersalne w architekturze klient\sphinxhyphen{}serwer.

\sphinxAtStartPar
Poniżej przedstawiono fragment skryptu aplikacyjnego odpowiedzialnego za odczyt z pliku i zasilenie tabel:

\begin{sphinxVerbatim}[commandchars=\\\{\}]
\PYG{k+kn}{import} \PYG{n+nn}{csv}
\PYG{k+kn}{import} \PYG{n+nn}{psycopg}

\PYG{c+c1}{\PYGZsh{} 1. Wczytywanie danych z pliku CSV do struktury iterowalnej (listy krotek)}
\PYG{k}{def} \PYG{n+nf}{wczytaj\PYGZus{}dane\PYGZus{}csv}\PYG{p}{(}\PYG{n}{sciezka\PYGZus{}pliku}\PYG{p}{)}\PYG{p}{:}
    \PYG{n}{dane} \PYG{o}{=} \PYG{p}{[}\PYG{p}{]}
    \PYG{k}{with} \PYG{n+nb}{open}\PYG{p}{(}\PYG{n}{sciezka\PYGZus{}pliku}\PYG{p}{,} \PYG{n}{mode}\PYG{o}{=}\PYG{l+s+s1}{\PYGZsq{}}\PYG{l+s+s1}{r}\PYG{l+s+s1}{\PYGZsq{}}\PYG{p}{,} \PYG{n}{encoding}\PYG{o}{=}\PYG{l+s+s1}{\PYGZsq{}}\PYG{l+s+s1}{utf\PYGZhy{}8}\PYG{l+s+s1}{\PYGZsq{}}\PYG{p}{)} \PYG{k}{as} \PYG{n}{plik}\PYG{p}{:}
        \PYG{n}{csv\PYGZus{}reader} \PYG{o}{=} \PYG{n}{csv}\PYG{o}{.}\PYG{n}{reader}\PYG{p}{(}\PYG{n}{plik}\PYG{p}{)}
        \PYG{n+nb}{next}\PYG{p}{(}\PYG{n}{csv\PYGZus{}reader}\PYG{p}{)}  \PYG{c+c1}{\PYGZsh{} Pominięcie wiersza nagłówka}
        \PYG{k}{for} \PYG{n}{wiersz} \PYG{o+ow}{in} \PYG{n}{csv\PYGZus{}reader}\PYG{p}{:}
            \PYG{n}{dane}\PYG{o}{.}\PYG{n}{append}\PYG{p}{(}\PYG{n+nb}{tuple}\PYG{p}{(}\PYG{n}{wiersz}\PYG{p}{)}\PYG{p}{)}
    \PYG{k}{return} \PYG{n}{dane}

\PYG{c+c1}{\PYGZsh{} 2. Właściwy proces wsadowego importu do bazy PostgreSQL}
\PYG{k}{def} \PYG{n+nf}{importuj\PYGZus{}do\PYGZus{}postgres}\PYG{p}{(}\PYG{n}{kategorie\PYGZus{}do\PYGZus{}importu}\PYG{p}{,} \PYG{n}{db\PYGZus{}config}\PYG{p}{)}\PYG{p}{:}
    \PYG{k}{try}\PYG{p}{:}
        \PYG{c+c1}{\PYGZsh{} Użycie Context Managera dla bezpiecznego zarządzania transakcjami}
        \PYG{k}{with} \PYG{n}{psycopg}\PYG{o}{.}\PYG{n}{connect}\PYG{p}{(}\PYG{o}{*}\PYG{o}{*}\PYG{n}{db\PYGZus{}config}\PYG{p}{)} \PYG{k}{as} \PYG{n}{conn}\PYG{p}{:}
            \PYG{k}{with} \PYG{n}{conn}\PYG{o}{.}\PYG{n}{cursor}\PYG{p}{(}\PYG{p}{)} \PYG{k}{as} \PYG{n}{cursor}\PYG{p}{:}

                \PYG{c+c1}{\PYGZsh{} Definicja sparametryzowanego zapytania zapobiegającego SQL Injection}
                \PYG{n}{zapytanie\PYGZus{}sql} \PYG{o}{=} \PYG{l+s+s2}{\PYGZdq{}}\PYG{l+s+s2}{INSERT INTO Kategorie (Nazwa\PYGZus{}Kategorii) VALUES (}\PYG{l+s+si}{\PYGZpc{}s}\PYG{l+s+s2}{)}\PYG{l+s+s2}{\PYGZdq{}}

                \PYG{c+c1}{\PYGZsh{} Wsadowe wywołanie zapytań (batch insert)}
                \PYG{n}{cursor}\PYG{o}{.}\PYG{n}{executemany}\PYG{p}{(}\PYG{n}{zapytanie\PYGZus{}sql}\PYG{p}{,} \PYG{n}{kategorie\PYGZus{}do\PYGZus{}importu}\PYG{p}{)}

                \PYG{n+nb}{print}\PYG{p}{(}\PYG{l+s+sa}{f}\PYG{l+s+s2}{\PYGZdq{}}\PYG{l+s+s2}{Pomyślnie zaimportowano }\PYG{l+s+si}{\PYGZob{}}\PYG{n+nb}{len}\PYG{p}{(}\PYG{n}{kategorie\PYGZus{}do\PYGZus{}importu}\PYG{p}{)}\PYG{l+s+si}{\PYGZcb{}}\PYG{l+s+s2}{ rekordów.}\PYG{l+s+s2}{\PYGZdq{}}\PYG{p}{)}

    \PYG{k}{except} \PYG{n+ne}{Exception} \PYG{k}{as} \PYG{n}{e}\PYG{p}{:}
        \PYG{n+nb}{print}\PYG{p}{(}\PYG{l+s+sa}{f}\PYG{l+s+s2}{\PYGZdq{}}\PYG{l+s+s2}{Wystąpił błąd operacji wsadowej: }\PYG{l+s+si}{\PYGZob{}}\PYG{n}{e}\PYG{l+s+si}{\PYGZcb{}}\PYG{l+s+s2}{\PYGZdq{}}\PYG{p}{)}
\end{sphinxVerbatim}


\section{Podsumowanie}
\label{\detokenize{rozdzial_4/index:podsumowanie}}
\sphinxAtStartPar
Niniejszy rozdział wieńczy etap projektowania i wdrażania struktury relacyjnej bazy danych dla systemu zarządzania biblioteką. Poprzez transformację modelu logicznego do postaci fizycznej, udało się z powodzeniem zaimplementować docelowe schematy w dwóch niezależnych środowiskach bazodanowych: produkcyjnym (PostgreSQL) oraz lokalnym/testowym (SQLite). Opracowane skrypty DDL poprawnie uwzględniają specyfikę obu silników, zachowując przy tym pożądaną architekturę relacyjną.

\sphinxAtStartPar
Dodatkowo, proces inicjalizacji systemu został zautomatyzowany dzięki przygotowaniu skryptu w języku Python. Wykorzystanie plików CSV oraz mechanizmu wsadowego wstawiania danych (batch insert) udowodniło, że możliwe jest szybkie i skalowalne zasilenie bazy niezbędnymi danymi demonstracyjnymi. Wdrożone rozwiązania stanowią solidny, zoptymalizowany fundament do dalszych prac nad logiką biznesową i warstwą aplikacyjną całego projektu.



\renewcommand{\indexname}{Index}
\printindex
\end{document}